% !TEX program = xelatex
\documentclass[12pt]{article}
\usepackage[margin=28mm]{geometry}
\usepackage{hyperref}
\usepackage{amsmath, amsthm, amssymb, amsfonts}
\usepackage{tikz-cd}
\usepackage{fontspec}
\usepackage{unicode-math}
\usepackage{xeCJK}
\setmainfont{Latin Modern Roman}
\setmonofont{Latin Modern Mono}
\setCJKmainfont{Noto Serif CJK JP}
\usepackage{listings}
\lstset{basicstyle=\ttfamily\footnotesize, columns=fullflexible, breaklines=true, showstringspaces=false}

\title{lean\_dual\_space\_task — 日本語ノート}
\author{TeX generated by ChatGPT-5 Pro\\Lean source generated by CODEX (OpenAI)}
\date{2025-09-10}

\begin{document}
\maketitle

\begin{abstract}
Bourbaki の精神（順序対と射影）:
  双対空間 `Module.Dual R V` は「射（線形汎関数）の集合」であり、
  等式は射の振る舞い（加法・スカラー倍）を読むことで成分的に確認できる。
  ここでは双対の和・評価・二重双対の評価、ペアリングの双線形性を
  既存の標準補題だけで示す。
\end{abstract}

\section*{生成元の明示（Provenance)}
\begin{itemize}
  \item \textbf{TeX 生成}: ChatGPT-5 Pro
  \item \textbf{Lean ソース生成}: CODEX (OpenAI)
\end{itemize}

\section*{概説}
本ノートはアップロードされた Lean ソースから、定義・定理の先頭行を抽出して表にまとめ、
併せて一般的な可換図式（課題名に応じて合成・核と像・双対空間の評価写像など）を掲示します。
著者名（author）には「TeXは ChatGPT-5 Pro により生成」「Lean ファイルは CODEX (OpenAI) により生成」を明記しています。

\section*{サマリー}
Total lines: 67\\
Defs: 0\\
Lemmas: 0\\
Theorems: 4\\
Structures: 0\\
Inductives: 0

\section*{図式（例）}

\[
\begin{tikzcd}
A \arrow[r, "f"] \arrow[d, "g"'] & B \arrow[d, "h"] \\
C \arrow[r, "k"'] & D
\end{tikzcd}
\]



\[
\begin{tikzcd}
V \times W^\* \arrow[r, "{\mathrm{id}\times f^\*}"] \arrow[dr, swap, "\langle -, -\rangle"] & V \times V^\* \arrow[d, "\langle -, -\rangle"] \\
& \Bbbk
\end{tikzcd}
\]


\section*{定理・定義の一覧（先頭行）}
\begin{center}
\begin{tabular}{lll}
\textbf{種別} & \textbf{名前} & \textbf{シグネチャ} \\ \hline
theorem & \texttt{dual\_add\_is\_linear} & \texttt{{R V : Type*} [CommSemiring R] [AddCommMonoid V] [Module R V] (φ ψ : Module.Dual R V) (v : V) : (φ + ψ) v = φ v + ψ v := by} \\
theorem & \texttt{eval\_map\_is\_linear} & \texttt{{R V : Type*} [CommSemiring R] [AddCommMonoid V] [Module R V] (v : V) (φ ψ : Module.Dual R V) : φ v + ψ v = (φ + ψ) v := by} \\
theorem & \texttt{double\_dual\_eval} & \texttt{{R V : Type*} [CommSemiring R] [AddCommMonoid V] [Module R V] (v : V) (φ : Module.Dual R V) : Module.Dual.eval R V v φ = φ v := by} \\
theorem & \texttt{dual\_pairing\_bilinear} & \texttt{{R V : Type*} [CommSemiring R] [AddCommMonoid V] [Module R V] (v w : V) (φ : Module.Dual R V) (a : R) : φ (a • v + w) = a * (φ v) + φ w := by} \\
\end{tabular}
\end{center}

\section*{Lean ソース（全文）}
\begin{lstlisting}
import Mathlib.Algebra.Module.LinearMap.Defs
import Mathlib.Algebra.Module.LinearMap.Basic
import Mathlib.LinearAlgebra.Dual.Defs

/-!
  Bourbaki の精神（順序対と射影）:
  双対空間 `Module.Dual R V` は「射（線形汎関数）の集合」であり、
  等式は射の振る舞い（加法・スカラー倍）を読むことで成分的に確認できる。
  ここでは双対の和・評価・二重双対の評価、ペアリングの双線形性を
  既存の標準補題だけで示す。
-/

noncomputable section

namespace MyProjects.LinearAlgebra.A

-- 双対空間の元（線形汎関数）の和の定義が線形であることの証明
theorem dual_add_is_linear {R V : Type*}
  [CommSemiring R] [AddCommMonoid V] [Module R V]
  (φ ψ : Module.Dual R V) (v : V) :
  (φ + ψ) v = φ v + ψ v := by
  -- これは線形写像の和の定義（点ごとの加法）。
  simpa [LinearMap.add_apply]

-- 評価写像の線形性：eval_v : V* → R が線形写像であること
theorem eval_map_is_linear {R V : Type*}
  [CommSemiring R] [AddCommMonoid V] [Module R V]
  (v : V) (φ ψ : Module.Dual R V) :
  φ v + ψ v = (φ + ψ) v := by
  -- 上の等式の対称形。右辺を定義で展開すれば自明。
  simpa [LinearMap.add_apply]

-- 発展：双対の双対への自然な埋め込み（有限次元では同型）
theorem double_dual_eval {R V : Type*}
  [CommSemiring R] [AddCommMonoid V] [Module R V]
  (v : V) (φ : Module.Dual R V) :
  Module.Dual.eval R V v φ = φ v := by
  -- 既存の単純化補題
  simpa using (Module.Dual.eval_apply (R:=R) (M:=V) v φ)

-- Bourbakiスタイル：双対ペアリング ⟨v, φ⟩ = φ(v) の双線形性
theorem dual_pairing_bilinear {R V : Type*}
  [CommSemiring R] [AddCommMonoid V] [Module R V]
  (v w : V) (φ : Module.Dual R V) (a : R) :
  φ (a • v + w) = a * (φ v) + φ w := by
  -- 射の加法性と斉次性を順に読む。`R` 上の `•` は乗法に一致。
  calc
    φ (a • v + w) = φ (a • v) + φ w := by simpa using (LinearMap.map_add φ (a • v) w)
    _ = a • φ v + φ w := by simpa using (LinearMap.map_smul φ a v)
    _ = a * (φ v) + φ w := by simpa [smul_eq_mul]

/-! ## Bourbaki スタイルの軽い example
  射（双対の元）の和・評価・二重双対評価を、定義と既存補題により確認。 -/

section Examples
  variable {R V : Type*}
  variable [CommSemiring R] [AddCommMonoid V] [Module R V]
  variable (φ ψ : Module.Dual R V) (v w : V) (a : R)

  example : (φ + ψ) v = φ v + ψ v := by simpa [LinearMap.add_apply]
  example : φ v + ψ v = (φ + ψ) v := by simpa [LinearMap.add_apply]
  example : Module.Dual.eval R V v φ = φ v := by simpa using (Module.Dual.eval_apply (R:=R) (M:=V) v φ)
  example : φ (a • v + w) = a * φ v + φ w := by
    simpa [smul_eq_mul] using (dual_pairing_bilinear (R:=R) (v:=v) (w:=w) (φ:=φ) (a:=a))
end Examples

end MyProjects.LinearAlgebra.A

\end{lstlisting}

\end{document}
